\subsection{The sharp four-site bound}

\begin{theorem}[Sharp four-site ratio]\label{thm:four-site-sharp}
For every four-element $H\subset\mathbb Z$ and all nonzero finitely
supported $f,g\ge0$,
\begin{equation}\label{eq:four-site-sharp}
 R_H(f,g)=\frac{\|f\star\widetilde{g1_H}\|_1}{\|f\star g\|_1}
 \le\frac74.
\end{equation}
The denominator retains all values of $g$, including those outside $H$.
The supremum over the inputs and the four integer sites is $K_4=7/4$.
\end{theorem}

The upper bound follows from an inequality for four events and finite
linear-programming duality. We first prove the probability inequality.

\begin{lemma}\label{lem:four-ratio-matrix}
Let $m=(a,b,c,d)$ with $a\ge b\ge c\ge d>0$ and $d/a\ge1/3$.
Define $r_{ij}=\min(m_i,m_j)/\max(m_i,m_j)$ for $0\le i,j\le3$,
and let $\mathcal R=(r_{ij})$. Then, as quadratic forms,
\[
 \mathcal R\succeq\frac58\,\mathbf1\mathbf1^{\mathrm T}.
\]
\end{lemma}
\begin{proof}
Set $q_1=b/a$, $q_2=c/b$, $q_3=d/c$, so $0<q_i\le1$ and
$q_1q_2q_3\ge1/3$.  For an orthonormal basis $e_0,e_1,e_2,e_3$,
define
\[
 v_0=e_0,\qquad
 v_i=q_i v_{i-1}+\sqrt{1-q_i^2}\,e_i.
\]
Their Gram matrix is $\mathcal R$.
Set
\[
 u=e_0+\sum_{i=1}^3\sqrt{\frac{1-q_i}{1+q_i}}\,e_i.
\]
Induction gives $u\cdot v_i=1$ for every $i$, including when
$q_i=1$.  If $t_i=(1-q_i)/(1+q_i)$ and $S=t_1+t_2+t_3$, then
$\|u\|^2=1+S$.

To bound $S$, put $P=t_1t_2+t_1t_3+t_2t_3$ and $T=t_1t_2t_3$.
The product condition gives
\[
 \begin{aligned}
 3(1-S+P-T)&\ge1+S+P+T,\\
 2S&\le1+P-2T\le1+\frac{S^2}{3}.
 \end{aligned}
\]
Thus $S^2-6S+3\ge0$.  As $0\le S<3$, this forces $S<3/5$:
on $[3/5,3]$ that polynomial is decreasing and its value at
$3/5$ is $-6/25$.  Hence $\|u\|^2<8/5$.
For any real vector $\xi$,
\[
 \left(\sum_i\xi_i\right)^2
 =\left(u\cdot\sum_i\xi_i v_i\right)^2
 \le\frac85\left\|\sum_i\xi_i v_i\right\|^2
 =\frac85\,\xi^{\mathrm T}\mathcal R\xi.
\]
This proves the claimed non-strict bound.
\end{proof}

\begin{lemma}[Four-event star inequality]\label{lem:four-event-star}
Let $J$ be a random subset of $\{0,1,2,3\}$, with
\[
 m_i=\mathbb P(i\in J)>0,\qquad
 c_{ij}=\mathbb P(i,j\in J),\qquad
 r_{ij}=\frac{\min(m_i,m_j)}{\max(m_i,m_j)}.
\]
Put
\[
 T=\sum_i m_i,\qquad
 D_i=\sum_{j\ne i}r_{ij}c_{ij},\qquad W=\max_iD_i.
\]
Then
\[
 T-W\le\frac74.
\]
Consequently the same bound holds when $W$ is replaced by the maximum
weight of a spanning tree, with edge weights $r_{ij}c_{ij}$.
The constant $7/4$ is attained for both versions.
\end{lemma}
\begin{proof}
\medskip\noindent\textbf{Reductions.}

We may relabel so that
\[
 1\ge a=m_0\ge b=m_1\ge c=m_2\ge d=m_3>0.
\]
The union bound gives $c_{0j}\ge a+m_j-1$, even when the right side
is negative.  Therefore
\begin{equation}\label{eq:four-star-frechet}
 T-D_0
 \le a+\frac{b(1-b)+c(1-c)+d(1-d)}a
 \le a+\frac{3}{4a}.
\end{equation}
For $3/4\le a\le1$ the last quantity is at most $7/4$, since
\[
 a+\frac{3}{4a}-\frac74
 =\frac{(4a-3)(a-1)}{4a}\le0.
\]

For any three positive-marginal events, their analogous maximum-star
deficit is at most $3/2$.  Indeed, if their largest marginal $u$ is
less than $1/2$, their total marginal is less than $3/2$.  Otherwise
the same union-bound argument gives a deficit at most
\[
 u+\frac1{2u}\le\frac32,\qquad \frac12\le u\le1,
\]
because the difference is $(2u-1)(u-1)/(2u)$.
Deleting the event with marginal $d$ and retaining a maximizing
three-event star thus proves the claimed bound when $d\le1/4$.

Let
\[
 A=\sum_{i<j}r_{ij}c_{ij},\qquad N=|J|.
\]
Always $W\ge A/2$, by averaging the four star degrees.
If $d>1/4$ and $a\le1/2$, all ratios are at least $1/2$, so
\[
 A\ge\frac12\mathbb E\binom N2\ge\frac{T-1}{2}.
\]
The second inequality follows from
$\binom k2\ge k-1$ for integers $0\le k\le4$.
As $T\le2$, we obtain
\[
 T-W\le T-\frac{T-1}{4}
       =\frac{3T+1}{4}\le\frac74.
\]
The case $T\le7/4$ is also immediate from nonnegative edge weights.

It remains to consider
\begin{equation}\label{eq:four-star-remaining}
 \frac12<a<\frac34,\qquad d>\frac14,\qquad T>\frac74.
\end{equation}
If the first bound in \eqref{eq:four-star-frechet} is at most $7/4$, there is
nothing more to prove.  Otherwise, write
\[
 a=\frac12+x,\quad b=\frac12+y,\quad c=\frac12+z,\quad
 d=\frac12+w,\quad s=y+z+w,\quad E=y^2+z^2+w^2.
\]
The failure of that bound implies
\begin{equation}\label{eq:four-star-ellipse}
 0<x<\frac14,\qquad
 E<R_0:=\left(\frac12-x\right)\left(\frac14-x\right)
       =x^2-\frac34x+\frac18.
\end{equation}
Indeed the numerator $b(1-b)+c(1-c)+d(1-d)$ is $3/4-E$.

\medskip\noindent\textbf{A variance estimate.}

Let $Y_i=1_{\{i\in J\}}$ and $m=(a,b,c,d)^{\mathrm T}$.
Under \eqref{eq:four-star-remaining}, $d/a>1/3$, so Lemma~\ref{lem:four-ratio-matrix} applies to
$Y-m$ pointwise.  Taking expectations yields
\[
 \mathbb E[Y^{\mathrm T}\mathcal RY]
 \ge m^{\mathrm T}\mathcal Rm+\frac58\operatorname{Var}N.
\]
The left side is $T+2A$, and the ordering of the marginals gives
\[
 m^{\mathrm T}\mathcal Rm=a^2+3b^2+5c^2+7d^2.
\]
Consequently
\begin{equation}\label{eq:four-star-variance-bound}
 T-W\le
 \frac{5T-(a^2+3b^2+5c^2+7d^2)-(5/8)\operatorname{Var}N}{4}.
\end{equation}
It is enough to prove
\begin{equation}\label{eq:four-star-target}
 D+\frac58\operatorname{Var}N\ge0,
 \qquad
 D=a^2+3b^2+5c^2+7d^2-5T+7.
\end{equation}
In the variables of \eqref{eq:four-star-ellipse}, $T=2+x+s$ and
direct expansion gives
\begin{equation}\label{eq:four-star-D}
 D=1-4x+x^2-2y+2w+3y^2+5z^2+7w^2.
\end{equation}

\medskip\noindent\textbf{The case $T\le2$.}

Here $s\le-x$.  We prove the stronger assertion $D>0$.
Completing squares and applying weighted Cauchy--Schwarz gives
\[
 \begin{split}
 D&=1-4x+x^2-\frac{10}{21}
       +3(y-1/3)^2+5z^2+7(w+1/7)^2\\
 &\ge1-4x+x^2-\frac{10}{21}
       +\frac{105}{71}(s-4/21)^2\\
 &\ge\frac{176x^2-244x+41}{71}.
 \end{split}
\]
For $0<x\le3/16$, the numerator is decreasing and is at least
$23/16$, so this is positive.

Suppose instead $x\ge3/16$.  The inequalities
$E\ge s^2/3\ge x^2/3$ and $E<R_0$ imply
\[
 R_0-\frac{x^2}{3}
 =\frac23x^2-\frac34x+\frac18>0.
\]
This quadratic is decreasing on $[0,1/4]$ and has value
$-1/1936$ at $9/44$, so $x<9/44$.

Orthogonal projection onto the plane perpendicular to $(1,1,1)$
gives
\[
 (y-w)^2\le2(E-s^2/3)\le2(E-x^2/3).
\]
Using $3y^2+5z^2+7w^2\ge3E$ in \eqref{eq:four-star-D}, we obtain
\[
 D\ge1-4x+x^2+3E-2\sqrt{2(E-x^2/3)}.
\]
As a function of $E$ on the allowed interval, the last two terms
are decreasing: their derivative is
$3-2/\sqrt{2(E-x^2/3)}<0$, since $E\le R_0\le1/8$.
(The lower endpoint follows by continuity.)
Thus
\[
 D\ge\frac{11}{8}-\frac{25x}{4}+4x^2
       -2\sqrt{2(R_0-x^2/3)}.
\]
For $3/16\le x<9/44$, the decreasing function
$R_0-x^2/3$ is at most $1/128$, its value at $3/16$.
Also $11/8-25x/4+4x^2$ is decreasing here and is at least
$511/1936$, its value at $9/44$.  Hence
\[
 D\ge\frac{511}{1936}-\frac14=\frac{27}{1936}>0.
\]
This proves \eqref{eq:four-star-target} throughout the case $T\le2$.

\medskip\noindent\textbf{The case $T\ge2$.}

Since $a<3/4$, we have $2\le T<3$.
For every integer $N$ the product $(N-2)(N-3)$ is nonnegative,
so
\[
 \operatorname{Var}N\ge(T-2)(3-T).
\]
Put $t=x+s=T-2$, $\rho=5/8$, and $h=\rho(1-2x)$.
By expansion,
\[
 \begin{split}
 D+\rho t(1-t)
 &=B+(-2+h)y+hz+(2+h)w\\
 &\quad+3y^2+5z^2+7w^2-\rho s^2,\\
 B&=1-\frac{27x}{8}+\frac{3x^2}{8}.
 \end{split}
\]
Weighted Cauchy--Schwarz gives
\[
 s^2\le\frac{23}{15}(y^2+3z^2+5w^2),
 \qquad \rho\,\frac{23}{15}=\frac{23}{24}<1.
\]
It follows that
\[
 3y^2+5z^2+7w^2-\rho s^2\ge2E.
\]
The squared length of the linear coefficient vector is
\[
 L^2=(-2+h)^2+h^2+(2+h)^2
     =8+\frac{75}{64}(1-2x)^2.
\]
Therefore
\[
 \begin{aligned}
 D+\rho t(1-t)&\ge B-L\sqrt E+2E\\
              &\ge B-L\sqrt{R_0}+2R_0\\
              &=Z-L\sqrt{R_0},
 \end{aligned}
\]
where $Z=(10-39x+19x^2)/8$.
The function $2E-L\sqrt E$ is decreasing for $0\le E\le R_0\le1/8$,
as follows either by differentiation or by setting $E=v^2$
and noting $4v-L<0$.  Also $Z\ge23/128>0$ for
$0\le x\le1/4$.

For completeness, the remaining rational polynomial check is
displayed explicitly:
\begin{equation}\label{eq:four-star-quartic}
 \begin{aligned}
 Z^2-L^2R_0&=\frac{F(4x)}{16384},\\
 F(u)&=61u^4-3828u^3+16824u^2-19344u+6816.
 \end{aligned}
\end{equation}
Let $B_{4,k}(v)=\binom4k v^k(1-v)^{4-k}$.
Direct expansion gives the following Bernstein coefficients:
\[
 \begin{array}{c|rrrrr}
  &k=0&k=1&k=2&k=3&k=4\\ \hline
 F(v/2)&6816&4398&2681&12363/8&14005/16\\
 F((1+v)/2)&14005/16&821/4&3/4&146&529
 \end{array}
\]
The table means that each polynomial in the left column equals
the sum of the listed coefficient times $B_{4,k}(v)$.
All coefficients are positive.  As $B_{4,k}(v)\ge0$ and
$\sum_kB_{4,k}(v)=1$ on $[0,1]$, we have $F(u)>0$ on $[0,1]$.
Equation \eqref{eq:four-star-quartic}, together with $Z>0$, gives
$Z>L\sqrt{R_0}$.  This proves \eqref{eq:four-star-target} for $T\ge2$.

Together the reductions and the two cases prove the star inequality.
Every star is a spanning tree, so the tree inequality follows.

\medskip\noindent\textbf{Equality.}

Give probability $1/6$ to each of
\[
 \{0,1\},\quad\{0,2\},\quad\{0,3\},\quad
 \{0,1,2\},\quad\{0,1,3\},\quad\{0,2,3\}.
\]
Then $m_0=1$ and $m_1=m_2=m_3=1/2$, while
$c_{0j}=1/2$ and $c_{ij}=1/6$ for distinct leaves.
Thus every center--leaf edge has weight $1/4$, and every
leaf--leaf edge has weight $1/6$.
The center star has weight $3/4$; each leaf star has weight
$7/12$.  Every spanning tree has three edges, each of weight
at most $1/4$, and the center star attains that bound.
As $T=5/2$, both deficits are exactly $7/4$.
\end{proof}


\begin{proof}[Proof of Theorem~\ref{thm:four-site-sharp}]
If $g$ vanishes at some point of $H$, remove its zero coordinates.
Corollary~\ref{cor:small-exact} bounds the resulting ratio by $3/2$;
if all coordinates vanish, the ratio is zero. Thus suppose
$H=\{p_0,p_1,p_2,p_3\}$ and $g_i=g(p_i)>0$ for every $i$.
Put $G=\sum_i g_i$ and $C=7/4$, and define
\[
 \rho_{ij}=\frac{\min(g_i,g_j)}{\max(g_i,g_j)},\qquad
 \kappa_{ij}=\frac{\min(g_i,g_j)}{g_i+g_j},
\]
\[
 \Phi_{ij}(v_i,v_j)=\min(v_i,v_j)
 +\min\left(\frac{g_j}{g_i}v_i,\frac{g_i}{g_j}v_j\right).
\]
Bounding both minima by their $j$ arguments gives
\[
 \kappa_{ij}\Phi_{ij}(v_i,v_j)
 \le\frac{\min(g_i,g_j)}{g_j}v_j\le v_j,
\]
and the analogous bound holds with $i$. For any star on the four
indices, root its edges at an index maximizing $v_i$ and charge each
edge to its child. Every other coordinate is charged once. Consequently,
if $w_i\ge0$ and $\sum_i w_i=1$ are probabilities for the four star
centers, then
\[
 \begin{aligned}
 W_w(v)&=\sum_{i<j}(w_i+w_j)\kappa_{ij}\Phi_{ij}(v_i,v_j),\\
 q_-(v)&=\sum_i v_i-W_w(v)\ge\max_i v_i.
 \end{aligned}
\]

We seek real diagonal coefficients $d_i$ with
\begin{equation}\label{eq:four-star-trace}
 \sum_i g_i d_i=G
\end{equation}
such that $q_+(v)=\sum_i d_i v_i-W_w(v)\le C\max_i v_i$.
The function $q_+$ is convex and positively homogeneous. Its upper
bound therefore follows from its values on the unit-cube vertices:
\begin{equation}\label{eq:four-star-vertices}
 \sum_{i\in J}d_i-
 \sum_{\{i,j\}\subset J}(w_i+w_j)\rho_{ij}\le C
 \qquad(\varnothing\ne J\subseteq\{0,1,2,3\}).
\end{equation}
Here $\kappa_{ij}\Phi_{ij}(1,1)=\rho_{ij}$; the empty vertex is zero.

Maximize $\sum_i g_i d_i$ subject to
\eqref{eq:four-star-vertices}, $w_i\ge0$ and $\sum_iw_i=1$.
This finite linear program is feasible with $d_i=0$.
The singleton constraints give $d_i\le C$. Since every $g_i>0$,
a maximizing level set can also be bounded below in every $d_i$,
so the maximum is finite and attained.
Its dual value is
\begin{equation}\label{eq:four-star-dual}
 \min_\lambda\left[
 C\sum_{J\ne\varnothing}\lambda_J+
 \max_i\sum_{j\ne i}\rho_{ij}
              \sum_{J\supset\{i,j\}}\lambda_J
 \right],
\end{equation}
where $\lambda_J\ge0$ is indexed by the nonempty subsets and
$\sum_{J\ni i}\lambda_J=g_i$ for every $i$.
Indeed, the multipliers of \eqref{eq:four-star-vertices} must have
the displayed marginals because the $d_i$ are unrestricted.
Their remaining objective is linear in $w$, so its maximum on the
probability simplex is the largest of the four star sums.

For each feasible $\lambda$, set $L=\sum_J\lambda_J>0$ and give the
nonempty subset $J$ probability $\lambda_J/L$. Its marginals are
$m_i=g_i/L$ and its pair marginals are
$c_{ij}=L^{-1}\sum_{J\supset\{i,j\}}\lambda_J$.
The ratios of the $m_i$ equal those of the $g_i$.
Lemma~\ref{lem:four-event-star} shows that the expression in
\eqref{eq:four-star-dual} is at least
$L\sum_i m_i=G$. Finite linear-programming duality therefore supplies
$w,d$ with objective at least $G$. Decreasing all $d_i$ by a common
nonnegative constant preserves \eqref{eq:four-star-vertices} and
gives the exact trace \eqref{eq:four-star-trace}.

Finally set $v_i^\pm(z)=g_i f(z\mp p_i)$.
Translation by $p_i+p_j$ interchanges the two terms of
$\Phi_{ij}(v_i^+,v_j^+)$ with those of
$\Phi_{ij}(v_i^-,v_j^-)$, so their sums over $z$ agree.
Also $\sum_zv_i^\pm(z)=g_i\|f\|_1$.
The trace condition and the two pointwise comparisons now give
\[
 \begin{aligned}
 \|f\star\widetilde{g1_H}\|_1
 &\le\sum_zq_-(v^-(z))
  =\sum_zq_+(v^+(z))\\
 &\le C\|f\star(g1_H)\|_1
 \le C\|f\star g\|_1.
 \end{aligned}
\]
This proves the upper bound with the full denominator.
Proposition~\ref{prop:geometric-product}, with four kernel values
equal to one, gives the matching lower bound on the supremum.
\end{proof}

\begin{corollary}[Sharp weighted bound on four sites]
\label{cor:four-site-weighted}
Let $k\ge0$ have at most four positive sites, and put
$K=\|k\|_1$ and $M=\|k\|_\infty$. For all nonzero finitely supported
$f,g\ge0$,
\begin{equation}\label{eq:four-site-weighted}
 \frac{\|f\star\widetilde{gk}\|_1}{\|f\star g\|_1}
 \le\frac{K+3M}{4}.
\end{equation}
For every prescribed nonempty list of at most four positive kernel
values, the supremum over the inputs and distinct integer locations
is exactly $(K+3M)/4$.
\end{corollary}
\begin{proof}
The zero kernel is immediate. Otherwise let
$F_t=\{p:k(p)>t\}$ for $0\le t<M$, and keep
$P=\|f\star g\|_1$ fixed. Every $F_t$ is nonempty and has at most
four elements. At each output, the maximum of the kernel integrals
is at most the integral of the maxima. Thus
Corollary~\ref{cor:small-exact} and
Theorem~\ref{thm:four-site-sharp} give
\[
 \begin{aligned}
 \|f\star\widetilde{gk}\|_1
 &\le\int_0^M\|f\star\widetilde{g1_{F_t}}\|_1\,dt\\
 &\le\frac P4\int_0^M(|F_t|+3)\,dt
  =\frac{K+3M}{4}\,P.
 \end{aligned}
\]
The sum over outputs is finite, and each application uses the same
profile and full denominator. Proposition~\ref{prop:geometric-product}
gives sharpness for each list of at least two values; a singleton
attains the value $M$. The lower construction permits the integer
locations to vary. It does not assert attainment on every fixed support.
\end{proof}
